\chapter{Distance, Norme et Produit Scalaire}

\justify

\setlength{\parindent}{0pt}
\renewcommand{\labelitemi}{\textbullet} % Utiliser des points noirs (•)

% ==================================================================================================================================
% Introduction

Plaçons nous dans un espace quelconque, par exemple des fonctions continues ou des matrice réelles de taille 4. 
Peut-on définir une notion de distance entre deux éléments de ces espaces ? Calculer leur norme ? 
Définir des sous-ensembles autrement que par des propriétés intrinsèques de ces objets ? 
Pour cela, nous allons avoir besoin de Normes, Produits Scalaires et Distances. 
Ce sont des applications particulières qui vont nous permettre de "mesurer" des "longueurs" dans ces espaces, 
et par la suite, construite de nouveau objets. 

\vspace{0.3cm}

Dans tout ce chapitre, sauf mention contraire, nous nous plaçons dans un $\R$-espace vectoriel, pas forcément en dimension finie. 


% ==================================================================================================================================
% Produit Scalaire

\section{Produit Scalaire}

Quand on parle de produit scalaire, on pense souvent à l'application dans le cas euclidien permettant de vérifier si 
deux vecteurs sont orthogonaux, en réalité, il existe une multitute de produits scalaires agissant sur tout autant d'espaces. 

\begin{definition}[Produit Scalaire]
    Soit $E$ un $\R$-espace vectoriel. Une application $\phi : E \times E \rightarrow \R$ de $E$ est un produit scalaire ssi c'est une forme :
    \begin{itemize}
        \item \textbf{Bilinéaire :} $ \forall x,y,z \in E, \forall \lambda,\mu \in \R$ on a :
            \begin{align}
                \phi(\lambda x + \mu y,z) = \lambda \phi(x,z) + \mu \phi(y,z) \\ 
                \phi(x,\lambda y + \mu z) = \lambda \phi(x,z) + \mu \phi(x,z)
            \end{align}
        \item \textbf{Symétrique :} $ \forall x,y \in E, \quad \phi(x,y) = \phi(y,x)$ 
        \item \textbf{Définie :} $ \forall x \in E, \quad \phi(x,x) = 0_\R \Longrightarrow x = 0_E $
        \item \textbf{Positive :} $ \forall x,y \in E, \quad \phi(x,y) \geq 0 $
    \end{itemize}
    On dit qu'un produit scalaire est une forme puisqu'elle est définie de E dans un corps telle que : 
        \[ \phi : E \longrightarrow \R \] 
    On note généralement un produit scalaire $ \langle .,. \rangle$ ou $(.|.)$. 
\end{definition}

\begin{remark}
    En pratique, pour montrer qu'une application est un produit scalaire, il suffit juste de montrer qu'elle est symétrique, 
    linéaire et définie positive. La bilinéarité découle de la symétrie et de la linéarité.  
\end{remark}

\begin{example}[Produits Sclaires]
    Regardons quelques exemples de produits scalaires...
    \begin{itemize}
        \item En génométrie euclidienne, on utilise un produit scalaire permettant de déterminer si deux vecteur sont 
        orthogonaux. Soient $A,B,C,D \in \R^2$, on définit alors :
            \[ \vec{AB}.\vec{CD} = AB \times CD \times \cos (\widehat{\vec{AB}\vec{CD}}) \]
        \item Dans l'espace des fonctions continues sur un intervalle $[a,b]$ on peut définir le produit scalaire suivant :
            \[ \forall f,g \in \mathcal{C}^0([a,b]), \quad \langle f,g \rangle = \int_{a}^{b} f(x)g(x) \; dx \] 
        \item Enfin, dans $\R^n$ on a le produit scalaire dit \textbf{euclidien} défini par : 
            \[ \forall x,y \in \R^n \text{ tels que : }
                \begin{cases}
                    x = (x_1, \dots, x_n) \\ 
                    y = (y_1, \dots, y_n)
                \end{cases}
                \text{ on a  : }
                \langle x,y \rangle = \sum_{i=1}^n (x_i y_i)^2 
            \] 
    \end{itemize}
\end{example}

\begin{prop}[Inégalité sur un produit scalaire]
    Soient $x,y \in E$, on a l'égalité suivante, valable pour tout produit scalaire :
        \[ \boxed{ \langle x,y \rangle ^2 \leq \langle x,x \rangle \times \langle y,y \rangle } \] 
\end{prop}

\begin{remark}
    Ce résultat découle de l'inégalité de Cauchy-Schwarz, vu plus tard. 
\end{remark}


% ==================================================================================================================================
% Norme


\section{Norme}

Une fois un produit scalaire défini sur un espace, on peut définir une application supplémentaire nous donnant plus d'informations 
sur un élément de l'espace. Nous allons donc définir une norme de deux façon, à partir d'un produit scalaire mais aussi 
de façon axiomatique. 

\begin{definition}[Norme]
    Soit $E$ un $\R$-espace vectoriel. Un application $\|.\| : E \rightarrow \R$ de $E$ est appelée \textbf{norme} ssi elle est : 
    \begin{itemize}
        \item \textbf{Définie : } $ \forall x \in E, \quad \|x \| = 0_\R \Longrightarrow x = 0_E $
        \item \textbf{Positive : } $ \forall x \in E, \quad \| x \| \geqslant 0 $
        \item \textbf{Positivement Homogène : } $ \forall x \in E, \forall \lambda \in \R, \quad \| \lambda x \| = |\lambda | \| x\| $  
        \item \textbf{Sous-additivité : } $ \forall x,y \in E, \quad \| x + y \| \leqslant \| x \| + \|y\| $
    \end{itemize}    
\end{definition}

\begin{remark}[Notation et vocabulaire]
    Tout comme le produit scalaire, une norme est une application vérifiant quelques propriétés. 
    Généralement, on note une norme $\|.\|$. Un espace muni d'une norme est appelé \textbf{espace normé}. 

    Dans un espace euclidien, une norme sert à "mesurer" la distance d'un point à l'origine. 
\end{remark}

\begin{example}
    Si on se place dans $\R$, le produit scalaire usuel sera la multiplication et la norme, la valeur absolue. 
\end{example}

\begin{proposition}[Norme à partir du produit scalaire]
    Dans un $\R$-espace vectoriel $E$, à partir d'un produit scalaire, on peut facilement définir une norme. 
    Soit $\langle .,. \rangle$ un produit scalaire sur $E$, alors l'application :
        \[
            \begin{cases}
                E \longrightarrow E \\
                x \longmapsto \sqrt{\langle x,x \rangle}
            \end{cases}
        \] 
    est une norme sur $E$. 
\end{proposition}

\begin{example}
    Dans $\R^n$, de même que le produit scalaire euclien, on peut définir la norme euclienne telle que :
    \[ 
        \forall x = (x_1, \dots, x_n) \in \R^n \quad \|x\| = \sqrt{ \sum_{i=1}^n (x_i)^2 }
    \] 
\end{example}

\begin{prop}[Inégalité de Cauchy-Schwarz]
    Soit $\langle .,. \rangle$ une produit scalaire sur $E$. On a alors l'inégalité suivante :
        \[ \boxed{ \forall x, y \in E, \quad | \langle x, y \rangle | \leqslant \| x\| \times \|y\| } \] 
    \emph{"La valeur absolue du produit scalaire est inférieure au produit des normes."}
\end{prop}


% ==================================================================================================================================
% Distance

\section{Distance}

Maintenant que nous pouvons "mesurer" des "longueurs" de vecteurs dans notre espace, on peut se demander si il est possible 
de "calculer" la distance entre deux éléments de $E$. Grâce à une telle application, on pourrait déterminer si deux éléments sont 
plus ou moins proches. Dès que l'on a une notion de distance, on peut ensuite l'intéresser à la notion de limite, etc... 

\begin{definition}[Distance]
    Soit $E$ un $\R$-espace vectoriel. 
    Une application $d(.,.) : E \times E \rightarrow \R$ est appelée \textbf{distance} ssi elle est :
    \begin{itemize}
        \item \textbf{Symétrique : } $ \forall x,y \in E, \quad d(x,y) = d(y,x) $
        \item \textbf{Définie : } $ \forall x,y \in E, \quad d(x,y) = 0_\R \Longrightarrow x = y $
        \item \textbf{Positive : } $ \forall x,y \in E, \quad d(x,y) \geqslant 0_\R $
        \item \textbf{Inégalité Triangulaire : } $ \forall x,yz \in E, \quad d(x,y) \leqslant d(x,z) + d(z,y)$ 
    \end{itemize}
\end{definition}

\begin{proposition}[Distance à partir d'une norme]
    Comme précédement, à partir d'une norme sur $E$, on peut facilement définir une distance entre deux vecteur. 
    Soit $\|.\|$ une norme sur $E$, l'application :
        \[ 
            \begin{cases}
                E \times E \longrightarrow \R \\ 
                (x,y) \longmapsto \| x - y \| 
            \end{cases}
        \] 
    Est une distance sur $E$. 
\end{proposition}

\begin{remark}[Vocabulaire]
    Un espace $E$ muni d'une distance est appelé \textbf{espace métrique}. 
\end{remark}


